\documentclass[a4paper,12pt,oneside]{article}
\makeatletter
\def\input@path{{../../}}
\makeatother
\usepackage{discrete-practice}

\begin{document}

\PracticeHeader{1}{4}{Бинарные отношения.}

\Task[Найдите область определения и множество значений отношений]{Задача 1}
\begin{alphlist}
  \item $\bigl\{(a,1),(a,2),(a,3),(a,4),(a,5),(a,6)\bigr\}$
  \item $\bigl\{(x,y)\,:\,x,y\in\mathbb{Z},\;x^2+y^2\le 16\bigr\}$
  \item $\bigl\{(x,y)\,:\,0\le x\le 10,\;0\le y\le 10,\;x>2y\bigr\}$
\end{alphlist}

\Task[Пусть отношения $U, V \subseteq \mathbb{R}\times\mathbb{R}$ определены указанным ниже способом]{Задача 2}
\[
U=\bigl\{(x,y)\,:\,y=x^2+5\bigr\},
\qquad
V=\bigl\{(x,y)\,:\,y=3x\bigr\}.
\]
\begin{alphlist}
  \item Опишите отношение $V\circ U$.
  \item Опишите отношение $U\circ V$.
  \item Опишите отношение $U^{-1}$.
  \item Опишите отношение $V^{-1}$.
  \item Найдите множество значений (образ) отношения $U$.
\end{alphlist}

\Task{Задача 3}
$A=\{a,b,c,d,e\}$. Опишите наименьшее рефлексивное отношение на этом множестве.

\Task{Задача 4}
\begin{numlist}
  \item Множество вершин $\{a,b,c,d,e,f\}$ и отношение $R$ для рёбер имеет вид:
  \[
    R_1=\bigl\{(a,b),(b,c),(a,c),(b,e),(c,f),(c,d),(d,f),(f,e)\bigr\}.
  \]

  \item Множество вершин $\{a,b,c,d\}$ и отношение $R$ для рёбер имеет вид:
  \[
    R_2=\bigl\{(b,c),(a,d),(b,a),(d,c),(b,d),(c,a),(d,d)\bigr\}.
  \]

  \item Определите свойства $R_1$ и $R_2$.
\end{numlist}

\Task{Задача 5}
Пусть $A=\{1,2,3\}$, а отношение $R$ на $A$ задано упорядоченными парами:
\[
R=\bigl\{(1,1),(2,1),(1,3),(3,1),(2,3)\bigr\}.
\]
Найдите рефлексивное, симметричное и транзитивное замыкания.

\Task{Задача 6}
Отношение $R_3$ задано на множестве $A=\{1,2,3,4,5\}$:
\[
a\,R_3\,b \iff a+b \text{ чётное}.
\]
Чему равна мощность множества классов эквивалентности $R_3$?

\Task{Задача 7}
Отношение $R$ на вещественной прямой $\mathbb{R}$ задано условием:
\[
xRy \iff x-y \text{ – целое число}.
\]
Докажите, что $R$ – отношение эквивалентности, и опишите классы эквивалентности,
содержащие $0$, $\tfrac{1}{2}$ и $\sqrt{2}$.

\end{document}
